\section{Polynomial expected depth with adaptive block modes}
\label{sec:adaptive-blocks}

Even after the first endpoint of a block answers positively, querying
the second endpoint can still help: it can pin the arc boundary and
reveal information that endpoint decoding would otherwise have to
recover. Crediting both effects improves the general expected-depth
factor to below $5.272565$. The proof separates the two tasks needed
for this bound. The local analysis below quantifies the saving from a
second endpoint, and a finite coverage-prefix certificate bounds the
cost of finding the block. Lemmas~\ref{lem:capped-rank-credit}
and~\ref{lem:stage-conversion} combine these local credits into the
recursive guarantee.

\subsection{Physical costs and a coarsened transcript}

Keep the polynomial exact-coverage and geometric-batch construction of
Section~\ref{sec:polynomial-blocks}. In a block of rank $k$, let $x,y$
be its newly hit masses exclusive to its two endpoints and $z$ their
shared mass. Orient the block so $x\le y$, and put $m=x+y+z$.
Each endpoint cap has mass at most $\delta w$, where $w=w(S)$;
thus $x+z,y+z\le\delta w$.

In \emph{stop mode}, stop querying this block at the first positive
answer. Relative to the nonentropy allowance $(2k+3)m$, the $y$ mass
exits second with a pinned boundary and saves one query per hypothesis.
The saving is $y$.

In \emph{query-both mode}, ask the second endpoint even when the first
was positive, querying the $y$ endpoint first. The coarsened event
$y+z$ then splits into its $y$ and $z$ parts when the second endpoint
is queried, which is the split charged by $E_*(y,z)$ below. Each
exclusive mass has one adjacent negative endpoint, so only one
alphabetic endpoint code is needed. Its nonentropy cost
is at most $2k+2$. The shared mass has two positive endpoints and
costs at most $2k+4$. Thus the total nonentropy cost is at most
\[
 (2k+2)(x+y)+(2k+4)z=(2k+2)m+2z.
\]
Subtracting this from the allowance $(2k+3)m$ leaves
$m-2z=x+y-z$, which is the physical saving.
An origin endpoint is free, and cached answers only lower these costs.
For the analysis, complete the chosen mode before singleton inference;
constant queries can be contracted afterwards.

Let $\mathcal R_{\rm act}$ be the actual discovery partition and
$J_{\rm act}=I(S,\mathcal R_{\rm act})$. Define a coarsened
partition $\mathcal R_*$ by forgetting the second answer on each
query-both route whose first answer was positive. This merges that
block's $y$ and $z$ events. Its $x$ event remains separate. Stop-mode
events are unchanged. Every nonzero coarsened event lies in a single
positive endpoint cap, so, for $J_*=I(S,\mathcal R_*)$,
\begin{equation}\label{eq:coarse-discovery-cap}
 J_*\ge b\sigma(\delta),\qquad b=w-p_0.
\end{equation}
The zero event is the same zero-pattern singleton as before.

The coarsening loses exactly
\[
 J_{\rm act}-J_*
 =\sum_{\text{query-both blocks}}
 (y+z)h\!\left(\frac{y}{y+z}\right),
\]
with zero contribution when $y+z=0$. Here
$h(p)=-p\ln p-(1-p)\ln(1-p)$. Actual endpoint decoding costs
at most its physical allowance plus $(I-J_{\rm act})/\ln2$,
where $I=I(S,\Pi)$ is the full-pattern information. Since
$J_{\rm act}-J_*$ is the sum displayed above, we can write
\[
 \frac{I-J_{\rm act}}{\ln2}
 =\frac{I-J_*}{\ln2}
   -\sum_{\text{query-both blocks}}
      (y+z)H_2\!\left(\frac{y}{y+z}\right).
\]
Here $H_2(p)=h(p)/\ln2$, so each term in the sum is the extra
endpoint information revealed by the second query. It therefore joins
the physical saving in the block ledger, while the remaining
information charge depends only on $J_*$. This accounts for the
second answer once, whether we view it as a reduction in decoding
cost or as information exposed during discovery.

\subsection{Rational mode choices and affine saving bounds}

Fix
\[
 p_* =\frac{21317}{125000},\qquad
 h_* =1-2p_* =\frac{41183}{62500}.
\]
Exact logarithm enclosures certify $H_2(p_*)>h_*$. Concavity and
symmetry of binary entropy give the rational tent bound
\[
 H_2(p)\ge\min\left\{\frac{h_*}{p_*}p,\ h_*,
                         \frac{h_*}{p_*}(1-p)\right\}.
\]
Its homogeneous mass version and the rational block saving are
\begin{equation}\label{eq:rational-block-saving}
 \begin{aligned}
 E_*(a,z)&=\min\left\{\frac{h_*}{p_*}a,\ h_*(a+z),
                                  \frac{h_*}{p_*}z\right\},\\
 D(x,y,z)&=\max\{y,\ x+y-z+E_*(y,z)\}.
 \end{aligned}
\end{equation}
Algorithm~\ref{alg:block-mode} selects the mode attaining this maximum;
Algorithm~\ref{alg:expected-tree} gives the full recursion. Both its choice and its saving
are rational functions of the supplied prior weights. Let
$\operatorname{Save}(G)=\sum_{\text{blocks of }G}D(x,y,z)$.
For each candidate geometric shift, choose its modes and minimize
the rational surrogate
\begin{equation}\label{eq:adaptive-shift-objective}
 F_*(G)=2\MSSC_{\rm blk}(G)+2p_0|G|-\operatorname{Save}(G).
\end{equation}
The full stage then obeys
\begin{equation}\label{eq:adaptive-stage}
 W_{\rm pat}(S)\le F_*(G)+3b+\frac{I-J_*}{\ln2}.
\end{equation}
Each mode preserves a block's union, newly covered mass, coverage sets,
and greedy order within its batch. The shift-average coverage proof
therefore remains available, and every shift satisfies
\eqref{eq:coarse-discovery-cap}.

\begin{algorithm}[H]
\caption{Rational choice between stopping and querying both endpoints}
\label{alg:block-mode}
\KwIn{New exclusive masses $x,y$, shared mass $z$;
calibrated $p_*$ and $h_*=1-2p_*$.}
\KwOut{Stored endpoint order, mode, and rational saving $D$.}
Label the endpoints so $x\le y$\;
$E\gets\min\{(h_*/p_*)y,\ h_*(y+z),\ (h_*/p_*)z\}$\;
\eIf{$y\ge x+y-z+E$}{
 Record stop mode with the $x$ endpoint first; $D\gets y$\;
}{
 Record query-both mode with the $y$ endpoint first;
 $D\gets x+y-z+E$\;
}
\KwRet the recorded order, mode, and $D$\;
\end{algorithm}

These new masses are computed on the all-negative prefix before the
block. In the stored discovery stage, stop mode ends at the first
positive answer; query-both mode learns both bits before endpoint
recovery. Their union and first-cover rank are the same.


\begin{lemma}[Affine block credits]\label{lem:affine-block-credit}
Put
\[
 c_* =\frac{228683}{103683}=\frac{2-p_*}{1-p_*},\qquad
 d_* =\frac{250000}{103683}=\frac2{1-p_*}.
\]
Every block with newly hit mass $m$ satisfies
\[
 D\ge0,\qquad D\ge m-\delta w,\qquad
 D\ge c_*m-d_*\delta w.
\]
\end{lemma}
\begin{proof}
The first two follow from $D\ge y\ge0$ and
$m-y=x+z\le\delta w$. If $\delta w=0$, these nonnegative masses
are all zero, so $m=0$ and the desired third bound is $D\ge0$.
We may therefore assume $\delta w>0$ and normalize $\delta w=1$.
If $m\le1$, then $c_*m-d_*\le c_*-d_*=-p_* /(1-p_*)<0$,
using $c_*>0$ and the displayed formulas for $c_*,d_*$. Otherwise
$1\le m\le2$, since $x+z\le1$ and $y+z\le1$ imply
$m=(x+z)+y\le2$. The same cap constraints give $x,y\ge m-1$,
while $z\ge0$ gives $x+y\le m$. Together with $x\le y$, these
inequalities define a triangle. Its vertices are the boundary
intersections $(x,y)=(m-1,m-1)$, $(m-1,1)$, and $(m/2,m/2)$:
\[
 (m-1,m-1),\qquad(m-1,1),\qquad(m/2,m/2).
\]
Each affine term defining $E_*(a,z)$ is nondecreasing in $a$, so
$E_*(y,z)\ge E_*(x,z)$ because $x\le y$. Thus the selected
query-both saving is at least $Q(x,y,z)=x+y-z+E_*(x,z)$. The function
$E_*(x,z)$ is the pointwise minimum of three affine functions, so it is
concave; adding the affine term $x+y-z$ preserves concavity of $Q$.
Every feasible point is a convex combination of its three
vertices, so concavity bounds its value below by the same convex
combination of the vertex values, and therefore by their minimum. The
two vertices with $z=0$ have value $m$. At the first vertex,
$x=y=m-1$ and $z=2-m$, so $x+y-z=3m-4$ and the three terms
in $E_*(x,z)$ are
\[
 E_*(m-1,2-m)=\min\left\{\frac{h_*}{p_*}(m-1),
 h_*\bigl((m-1)+(2-m)\bigr),\frac{h_*}{p_*}(2-m)\right\}
 =\min\left\{\frac{h_*}{p_*}(m-1),h_*,
 \frac{h_*}{p_*}(2-m)\right\}.
\]
Hence its value is
\[
 f(m)=3m-4+
       \min\left\{\frac{h_*}{p_*}(m-1),\ h_*,
                            \frac{h_*}{p_*}(2-m)\right\}.
\]
For $m\le1+p_*$, compare the normalized coverage credit $m-1$ with
the target line:
\[
 (m-1)-(c_*m-d_*)=(d_*-1)-(c_*-1)m
 =\frac{1+p_*-m}{1-p_*}\ge0.
\]
For every feasible point, $y=m-(x+z)\ge m-1$. Hence the stop-mode
saving $y$ dominates $c_*m-d_*$ throughout this range. On
$[1+p_*,2]$, the minimum in $f(m)$ is a concave function:
it is the pointwise minimum of three affine functions. Adding the
affine term $3m-4$ preserves concavity. At the left endpoint, the
first two entries in the minimum both equal $h_*$. The third is
$h_*(1-p_*)/p_*$, which is at least $h_*$ because $1-p_*>p_*$.
Thus $f(1+p_*)=3(1+p_*)-4+h_*=p_*$, using $h_*=1-2p_*$. At
the right endpoint, the third entry is zero and the other two are
positive, so $f(2)=3\cdot2-4=2$. The affine target line has the same
endpoint values:
\[
 c_*(1+p_*)-d_*=p_*=f(1+p_*),\qquad
 2c_*-d_*=2=f(2),
\]
where the first identity follows from
$(2-p_*)(1+p_*)-2=p_*(1-p_*)$, and the second from
$2(2-p_*)-2=2(1-p_*)$. Hence $c_*m-d_*$ is the chord joining the
endpoint values of $f$.
Concavity places $f$ on or above this chord. The other
vertex value $m$ also lies above the target line throughout $m\le2$,
because
\[
 m-(c_*m-d_*)=d_*-(c_*-1)m=\frac{2-m}{1-p_*}\ge0.
\]
Therefore all vertex values of the query-both expression dominate
$c_*m-d_*$ on the remaining range. Concavity then places its value at
every point of the triangle above this same target, so $D$ also exceeds
it. All terms scale linearly with the masses; multiplying by
$\delta w$ therefore proves the original bound.
\end{proof}

\subsection{A certified mixture of coverage-prefix credits}

Keep the witness masses $m_t$ and $M_k=\sum_{t\le k}m_t$ from
Section~\ref{sec:polynomial-blocks}, with $m_t\le2\delta w$.
A comparison edge sequence covers at least $M_k$ in its first $k$
blocks. Exact budgeted coverage therefore gives covered mass at least
$M_k$ in a batch with budget $k$. If the schedule prefix through
that batch uses at most $\ell$ blocks, Lemma~\ref{lem:affine-block-credit}
and nonnegative later savings imply each of the lower bounds
\[
 \operatorname{Save}(G)\ge
 \mathsf A(k,\ell):=M_k-\delta w\ell,
 \qquad
 \operatorname{Save}(G)\ge
 \mathsf C(k,\ell):=c_*M_k-d_*\delta w\ell.
\]
It is also at least zero. Early completion only improves these bounds.
For a shift $u\in[1,\beta)$, the first batch has
$k=\lfloor u\rfloor$, $\ell\le\lfloor u\rfloor$; through the
second batch, $k=\lfloor\beta u\rfloor$ and
$\ell\le\lfloor u\rfloor+\lfloor\beta u\rfloor$.
Any convex combination of their lower bounds and zero bounds the
\emph{same} total saving, counting it once.

Use
\[
 \beta=\frac{3512863}{10^6},\qquad
 \delta=\frac{144136826}{10^9},\qquad
 v_1=\frac{830027}{10^7},\quad
 v_2=\frac{472577}{10^6},\quad
 v_3=\frac{1701301}{5\cdot10^6}.
\]
The following fixed choices are used only in the proof:
\begin{center}
\small
\begin{tabular}{ll}
\hline
Shift range & Lower bound on $\operatorname{Save}(G_u)$\\
\hline
$[1,4/\beta)$ & $\mathsf C(1,1)$\\
$[4/\beta,5/\beta)$ & $v_1\mathsf C(1,1)+(1-v_1)\mathsf A(4,5)$\\
$[5/\beta,6/\beta)$ & $\mathsf A(5,6)$\\
$[6/\beta,2)$ & $\mathsf C(1,1)$\\
$[2,10/\beta)$ & $\mathsf A(2,2)$\\
$[10/\beta,3)$ & $v_2\mathsf A(2,2)+(1-v_2)0$\\
$[3,11/\beta)$ & $v_3\mathsf A(3,3)+(1-v_3)\mathsf C(3,3)$\\
$[11/\beta,12/\beta)$ & $\mathsf C(3,3)$\\
$[12/\beta,\beta)$ & $\mathsf A(3,3)$\\
\hline
\end{tabular}
\end{center}

Here is an explicit coefficient definition for the certificate. Start
with $a_t=e_\beta(t)$ and $\mu=0$. A table row with range $[l,r)$
has log-uniform probability $P=\ln(r/l)/\ln\beta$. For each
minorant $\alpha M_k-\rho\delta w$ in that row, with mixing weight
$v'$, add $Pv'\alpha$ to every $a_t$ with $t\le k$ and add
$Pv'\rho$ to $\mu$. For a zero term both additions are zero.
No table entry has $k>5$, so $a_t=e_\beta(t)$ for $t\ge6$.
These fixed coefficients certify the analysis; the algorithm itself
compares the finite set of candidate shifts and modes.
The shift-average transfer is
\[
 \sum_t m_t e_\beta(t)+\mathbb E[\operatorname{Save}(G_u)]
 \ge\sum_t m_t a_t-\delta w\mu.
\]

Set $\tau=345603/250000=1.382412$. Exact rational logarithm
enclosures certify that only $a_4$ lies below $\tau$:
\[
 \begin{aligned}
 a_1&>1.38241213,& a_2&>1.38241211,\\
 a_3&>1.38241211,& a_4&<1.37612333,\\
 a_5&>1.38241212,& e_\beta(6)&>1.38382029.
 \end{aligned}
\]
Monotonicity of $e_\beta$ controls the remaining tail. Define
\begin{equation}\label{eq:adaptive-chi-B}
 \chi_* =\tau-2\delta(\tau-a_4)-\delta\mu,\qquad
 B_* =3+2\Phi_\beta(1)-\chi_*
      =7-\frac{\ln6}{\ln\beta}-\chi_*.
\end{equation}
This is Lemma~\ref{lem:capped-rank-credit} with the single
exceptional index $E=\{4\}$ and $\zeta=2/\ln\beta$.
The certified inequality $\chi_*+2/\ln\beta>\tau$ is its
zero-route condition. Thus the combined discrete, block-mode, and
zero-route discounts save at least $\chi_*b$.
Minimizing \eqref{eq:adaptive-shift-objective} allows the saving to
be retained under the same shift average as the first-hit and length
terms. The proof of \eqref{eq:poly-block-discrete} now gives
\begin{equation}\label{eq:adaptive-local}
 W_{\rm pat}(S)\le A\Delta+B_*b+\frac{I-J_*}{\ln2},
 \qquad A=\frac{\beta+1}{\ln\beta}.
\end{equation}
The at-most-two-real-tests degenerate-hull construction also satisfies
this bound because $A,B_*>2$, by the same argument as in the earlier
block proof.

\begin{theorem}[Adaptive-mode polynomial expected-depth approximation]
\label{thm:poly-adaptive-expected}
For general supplied binary unit-cost planar line queries and arbitrary
strictly positive rational priors, the adaptive-mode block algorithm
runs in deterministic polynomial time and has approximation factor
\[
 \boxed{\rho_*=
 A+\max\{\ln2/h(\delta),\ B_*\ln2/\sigma(\delta)\}<5.272565.}
\]
Its exponent has no dependence on the smallest prior.
\end{theorem}
\begin{proof}
Apply Lemma~\ref{lem:stage-conversion} with $B=B_*$ to
\eqref{eq:adaptive-local} and \eqref{eq:coarse-discovery-cap}.
The coarsened transcript is a function of the full pattern, so
$0\le J_*\le I$ and the lemma applies unchanged.

The exact rational certificate \texttt{further-expected-check.py}
verifies the tent inequality, all coefficient comparisons, the
zero-route condition, and
\[
 \begin{aligned}
 0.856049452254&<\chi_*<0.856049452257,\\
 4.717880246117&<B_*<4.717880246121,\\
 \max\{\ln2/h(\delta),B_*\ln2/\sigma(\delta)\}
 &<1.680753968384.
 \end{aligned}
\]
Together with $A<3.591810178840040$, these certify
$\rho_*<5.272564147225<5.272565$.
The script also checks 5,000 rational block-credit instances and
10,000 analytical source profiles. Exact interval bounds certify the
theorem's numerical claims.

Coverage DPs, finite shift enumeration, and the disjoint-pattern
recursion remain polynomial. Each mode uses fixed linear forms in the
prior weights. Computing the entropy, logarithms, or
analysis-table mixture is unnecessary for constructing the tree.
Every query remains a supplied endpoint line.
\end{proof}

\paragraph{Endpoint orientation.}
Each affine form defining $E_*(a,z)$ is nondecreasing in $a$, so the
query-both mode obtains at least as much information credit by asking
the larger exclusive mass first. The coarsened exits still lie in the
same endpoint caps. Stop mode uses the smaller exclusive mass first,
which leaves the larger exclusive class as the second exit. The
algorithm compares these two modes by their rational scores. The
affine certificate is tight when $x=y$, so changing the endpoint order
alone cannot improve its constant.
